\documentclass[12pt]{article}

\usepackage[utf8]{inputenc}
\usepackage[T1]{fontenc}
\usepackage[brazil]{babel}
\usepackage{geometry}
\usepackage{graphicx}
\usepackage{booktabs}
\usepackage{url}
\usepackage{hyperref}

\geometry{a4paper, margin=2.5cm}

\title{Recuperacao de Artigos Cientificos sobre o Uso de Inteligencia Artificial na Educacao}
\author{Guilherme de Araujo Brandao}
\date{2026}

\begin{document}

\maketitle

\begin{abstract}
This paper presents an information retrieval system for scientific articles about the use of Artificial Intelligence in Education. The system compares BM25, TF-IDF/KNN, and a hybrid sparse-dense ranking strategy. The evaluation is based on manually created queries and relevance judgments, using P@10, R@10, MAP, and nDCG@10.
\end{abstract}

\begin{resumo}
Este trabalho apresenta um sistema de recuperacao de artigos cientificos sobre o uso de Inteligencia Artificial na Educacao. O sistema compara BM25, KNN com TF-IDF e uma estrategia de ranking hibrido. A avaliacao utiliza consultas e julgamentos de relevancia construidos manualmente, com as metricas P@10, R@10, MAP e nDCG@10.
\end{resumo}

\section{Introducao}

A producao cientifica sobre Inteligencia Artificial na Educacao tem crescido com o avanco de sistemas tutores inteligentes, learning analytics, modelos generativos e grandes modelos de linguagem. Nesse contexto, torna-se necessario construir mecanismos capazes de recuperar artigos relevantes a partir de consultas textuais.

Este trabalho implementa o componente de recuperacao do paradigma Retrieval-Augmented Generation (RAG), concentrando-se no retorno de uma lista ranqueada de artigos cientificos.

\section{Trabalhos Relacionados}

Descrever BM25, TF-IDF, KNN aplicado a recuperacao, recuperacao densa e sistemas RAG.

\section{Metodologia}

\subsection{Escopo da colecao}

O tema definido foi o uso de Inteligencia Artificial na Educacao. Foram utilizadas palavras-chave relacionadas a intelligent tutoring systems, learning analytics, educational data mining, generative AI in education, ChatGPT in education e personalized learning.

\subsection{Coleta dos dados}

A colecao foi construida a partir da API do ArXiv. Para cada artigo, foram armazenados id, titulo, resumo, autores, categorias, data e URL.

\subsection{Pre-processamento}

O texto de cada documento foi formado pela concatenacao do titulo com o resumo. Em seguida, foram aplicadas normalizacao para minusculas, remocao de pontuacao, tokenizacao e remocao de stopwords.

\subsection{Recuperador BM25}

O BM25 foi utilizado como baseline esparso. Os parametros adotados foram $k_1=1.5$ e $b=0.75$.

\subsection{Recuperador KNN com TF-IDF}

O segundo recuperador representa documentos e consultas por vetores TF-IDF. A similaridade entre consulta e documento foi calculada com similaridade do cosseno, retornando os documentos mais proximos.

\subsection{Ranking hibrido}

Como modulo de aprofundamento, foi implementado um ranking hibrido combinando os scores normalizados de BM25 e KNN/TF-IDF.

\section{Avaliacao Experimental}

Foram criadas consultas representativas do tema e julgamentos de relevancia no formato qrels. As metricas utilizadas foram P@10, R@10, MAP e nDCG@10.

\section{Resultados e Discussao}

Inserir tabela com resultados dos tres sistemas: BM25, KNN/TF-IDF e hibrido. Discutir em quais consultas cada modelo apresentou melhor desempenho.

\section{Conclusao}

O trabalho apresentou um pipeline completo de recuperacao de artigos cientificos sobre IA na Educacao. Os resultados permitiram comparar abordagens esparsas, vetoriais e hibridas.

\section{Uso de IA Generativa}

Foi utilizado assistente de IA generativa como apoio para organizacao do projeto, estruturacao inicial do codigo e revisao textual. As decisoes metodologicas, execucao dos experimentos e analise dos resultados foram conduzidas pelo autor.

\bibliographystyle{plain}
\begin{thebibliography}{9}

\bibitem{robertson2009bm25}
Robertson, S. and Zaragoza, H. The Probabilistic Relevance Framework: BM25 and Beyond. Foundations and Trends in Information Retrieval, 2009.

\bibitem{manning2008ir}
Manning, C. D., Raghavan, P. and Schutze, H. Introduction to Information Retrieval. Cambridge University Press, 2008.

\bibitem{lewis2020rag}
Lewis, P. et al. Retrieval-Augmented Generation for Knowledge-Intensive NLP Tasks. NeurIPS, 2020.

\end{thebibliography}

\end{document}
